% ============================================================================
%  Cover Letter — GAHIB
%
%  Journal retargeting
%  -------------------
%  Change ALL three \newcommand lines in the "JOURNAL PROFILE" block below
%  to retarget this letter to a different journal. The rest of the letter
%  (Contributions, Declarations, closing) is journal-agnostic.
%
%  Ready-to-swap examples for other suitable venues are given as commented
%  alternatives under the JOURNAL PROFILE block — uncomment one and comment
%  out the Frontiers in Genetics defaults.
% ============================================================================
\documentclass[11pt,a4paper]{article}

% ── Packages ──────────────────────────────────────────────────────────────
\usepackage[T1]{fontenc}
\usepackage{lmodern}
\usepackage{microtype}
\usepackage{xcolor}
\usepackage{hyperref}
\usepackage{graphicx}
\usepackage{geometry}
\usepackage{enumitem}
\usepackage{titlesec}
\usepackage{fancyhdr}
\usepackage{tikz}
\usepackage{parskip}

% ── Page geometry ─────────────────────────────────────────────────────────
\geometry{
  left=2.4cm, right=2.4cm,
  top=1.8cm, bottom=1.8cm,
}

% ============================================================================
%  JOURNAL PROFILE  —  edit these three macros to retarget the letter
% ============================================================================

%  (1) Full journal name as it appears on the masthead (italicised in body).
\newcommand{\journalFull}{Frontiers in Genetics}

%  (2) Short name / acronym used in closing lines.
\newcommand{\journalShort}{\emph{Frontiers in Genetics}}

%  (3) Per-journal scope statement + bulleted fit. Keep 3--5 bullets.
\newcommand{\journalScope}{%
The manuscript addresses topics central to \journalShort, and in
particular its \emph{Computational Genomics} specialty section:
\begin{itemize}[leftmargin=1.4em, itemsep=1pt, topsep=2pt]
  \item \textbf{Single-cell genomics \& methods} --- a new
    representation-learning framework for scRNA-seq that unifies
    graph attention, a 2D information bottleneck, and
    Lorentz-hyperbolic geometry within one variational objective.
  \item \textbf{Machine learning for genomics} --- 21 baseline methods
    benchmarked across 53 datasets and 20 metrics with paired
    Wilcoxon signed-rank testing and Benjamini--Hochberg correction.
  \item \textbf{Biological interpretability} --- unsupervised
    developmental hierarchy recovery, GO biological-process enrichment,
    and decoder-Jacobian gene attribution.
  \item \textbf{Reproducibility \& open tools} --- open-source library,
    deterministic seeding, and multi-seed robustness analysis.
\end{itemize}%
}

% ── Alternative journal profiles (uncomment one to retarget) ─────────────
%
% %  Bioinformatics (Oxford)
% \renewcommand{\journalFull}{Bioinformatics}
% \renewcommand{\journalShort}{\emph{Bioinformatics}}
% \renewcommand{\journalScope}{%
% This work fits squarely within \journalShort's scope:
% \begin{itemize}[leftmargin=1.4em, itemsep=1pt, topsep=2pt]
%   \item \textbf{Genome analysis \& single-cell methods} --- a new
%     representation-learning framework for scRNA-seq.
%   \item \textbf{Statistical learning} --- paired Wilcoxon signed-rank
%     testing with Benjamini--Hochberg correction across 53 datasets.
%   \item \textbf{Reproducibility} --- open-source library, deterministic
%     seed, and multi-seed robustness analysis.
% \end{itemize}%
% }
%
% %  Briefings in Bioinformatics (Oxford)
% \renewcommand{\journalFull}{Briefings in Bioinformatics}
% \renewcommand{\journalShort}{\emph{Briefings in Bioinformatics}}
% \renewcommand{\journalScope}{%
% The manuscript addresses topics central to \journalShort:
% \begin{itemize}[leftmargin=1.4em, itemsep=1pt, topsep=2pt]
%   \item \textbf{Deep learning for biological data} --- graph attention,
%     information bottleneck, and hyperbolic geometry unified for
%     scRNA-seq.
%   \item \textbf{Benchmarking \& method comparison} --- 11 studies,
%     20 metrics, 53 datasets, 21 baseline methods.
%   \item \textbf{Biological interpretability} --- tissue-specific GO
%     enrichment and unsupervised developmental hierarchy recovery.
% \end{itemize}%
% }
%
% %  NAR Genomics and Bioinformatics (Oxford)
% \renewcommand{\journalFull}{NAR Genomics and Bioinformatics}
% \renewcommand{\journalShort}{\emph{NARGAB}}
%
% %  IEEE/ACM Transactions on Computational Biology and Bioinformatics
% \renewcommand{\journalFull}{IEEE/ACM Transactions on Computational
%   Biology and Bioinformatics}
% \renewcommand{\journalShort}{\emph{TCBB}}
%
% %  Patterns (Cell Press)
% \renewcommand{\journalFull}{Patterns}
% \renewcommand{\journalShort}{\emph{Patterns}}
%
% %  Nature Communications / Nature Machine Intelligence
% \renewcommand{\journalFull}{Nature Machine Intelligence}
% \renewcommand{\journalShort}{\emph{Nature Machine Intelligence}}

% ============================================================================
%  END JOURNAL PROFILE
% ============================================================================

% ── Colours ───────────────────────────────────────────────────────────────
\definecolor{accentblue}{HTML}{1A5276}
\definecolor{lightgray}{HTML}{F2F2F2}
\definecolor{darkgray}{HTML}{4A4A4A}
\definecolor{ruleblue}{HTML}{2980B9}

% ── Hyperlinks ────────────────────────────────────────────────────────────
\hypersetup{
  colorlinks=true,
  linkcolor=accentblue,
  urlcolor=ruleblue,
  pdfauthor={Zeyu Fu},
  pdftitle={Cover Letter — GAHIB},
}

% ── Header / footer ──────────────────────────────────────────────────────
\pagestyle{fancy}
\fancyhf{}
\renewcommand{\headrulewidth}{0pt}
\fancyfoot[C]{\small\color{darkgray}%
  Zeyu Fu\enspace$\cdot$\enspace
  Army Medical University, Chongqing\enspace$\cdot$\enspace
  \href{mailto:fuzeyu99@126.com}{fuzeyu99@126.com}}

% ── Section styling ──────────────────────────────────────────────────────
\titleformat{\section}
  {\large\bfseries\color{accentblue}}{\thesection}{0.6em}{}
  [\vspace{-4pt}{\color{ruleblue}\titlerule[0.8pt]}]

% ── Document ──────────────────────────────────────────────────────────────
\begin{document}
\thispagestyle{fancy}

% ── Letterhead ────────────────────────────────────────────────────────────
\noindent
\begin{minipage}[t]{0.62\linewidth}
  {\Large\bfseries\color{accentblue} Zeyu Fu}\\[3pt]
  {\small\color{darkgray}
    Institute of Combined Injury, College of Preventive Medicine\\
    Army Medical University, Chongqing 400038, China\\
    \href{mailto:fuzeyu99@126.com}{fuzeyu99@126.com}\\
    ORCID: \href{https://orcid.org/0009-0001-8329-0108}{0009-0001-8329-0108}}
\end{minipage}%
\hfill
\begin{minipage}[t]{0.34\linewidth}
  \raggedleft
  {\small\color{darkgray}\today}
\end{minipage}

\vspace{4pt}
{\color{ruleblue}\hrule height 1.2pt}
\vspace{8pt}

% ── Addressee ─────────────────────────────────────────────────────────────
\noindent
{\small\color{darkgray}%
\textbf{To the Editor-in-Chief}\\
\textit{\journalFull}}

\vspace{6pt}

\noindent\textbf{\large\color{accentblue}%
  Re:\ Submission of Original Research Article ---\\[2pt]
  ``GAHIB: Graph Attention VAE with a Hyperbolic Information
  Bottleneck\\
  for Biologically Structured Single-Cell Representations''}

\vspace{8pt}

% ── Opening ───────────────────────────────────────────────────────────────
\noindent Dear Editor,

\vspace{2pt}

I am submitting the above-titled manuscript for consideration as
an original research article in \textit{\journalFull}.
The paper presents \textbf{GAHIB}, a deep generative framework that
integrates graph attention, information-bottleneck compression, and
Lorentz-hyperbolic geometry within a single variational autoencoder
for learning biologically structured latent representations from
single-cell RNA-sequencing data.

% ── Graphical abstract ────────────────────────────────────────────────────
\vspace{4pt}
\noindent
\begin{center}
\includegraphics[width=\textwidth]{graphical_abstract.pdf}
\end{center}
\vspace{-2pt}

{\small
\noindent\textbf{Figure overview.}\enspace
(a)~GAHIB architecture: a $k$NN graph attention encoder feeds an
information-bottleneck compressor that is mapped to a Lorentz manifold
and decoded by a negative-binomial head, trained end-to-end with a
four-term loss.
(b)~Evaluation scope: 53 datasets, 20 metrics across seven benchmark
studies, with Wilcoxon signed-rank significance testing.
(c)~Dataset collection: 27 cancer and 26 developmental scRNA-seq
datasets (2{,}531--957{,}975 cells) spanning diverse tissues and scales.\par}

% ── Contributions ─────────────────────────────────────────────────────────
\section*{Contributions}

GAHIB combines a graph-attention encoder, a 2D information bottleneck,
and a Lorentz-hyperbolic geometry loss within one variational
objective; an ablation across 53 datasets shows that removing any of
these components reduces part of the performance profile (paired
Wilcoxon signed-rank, $p<0.05$, BH-corrected).
Across 11 studies spanning 7 deep-learning, 5 classical DR,
5 geometric-VAE, and 4 disentanglement baselines on 20 metrics,
GAHIB performs competitively across these baseline families, with
BH-corrected improvements in many pairwise comparisons.
Lorentz norms correlate with developmental hierarchy ($\rho_s=0.329$
mean across 53 datasets, up to 0.673 in enriched contexts), and the
derived pseudotime tracks diffusion pseudotime.
Decoder-Jacobian analysis recovers 31.5\% of marker genes, and
GO biological-process enrichment shows individual latent dimensions
align with tissue-specific programmes without supervision.
Multi-seed runs show $<$3\% standard deviation across metrics; at
fixed 200 epochs, training time changes little with cell count under
fixed subgraph sampling.

% ── Fit to journal ────────────────────────────────────────────────────────
\section*{Relevance to \journalShort}

\journalScope

% ── Declarations ──────────────────────────────────────────────────────────
\section*{Declarations}

\begin{itemize}[leftmargin=1.4em, itemsep=1pt, topsep=2pt]
  \item Not published previously; not under consideration elsewhere.
  \item All authors have read and approved the submitted version.
  \item No conflicts of interest to disclose.
  \item Funding: Special Project for Enhancing the Scientific and
    Technological Innovation Capacity of the Army Medical University
    (No.~2023XJS34).
  \item Source code and trained models:
    \url{https://github.com/PeterPonyu/GAHIB}.
\end{itemize}

\vspace{4pt}

% ── Closing ───────────────────────────────────────────────────────────────
\noindent
We believe the manuscript would be of interest to \journalShort{}
readers and would appreciate your consideration. We look forward to
the editors' and reviewers' feedback.

\vspace{12pt}

\noindent
Sincerely,\\[10pt]
{\bfseries\color{accentblue} Zeyu Fu}\hspace{0.6em}%
{\small\color{darkgray}---\ Institute of Combined Injury,
  Army Medical University, Chongqing, China}

\end{document}
